\section{Realizing the sharp profiles in one measure}
\label{realization:section}

The matching profiles are not excluded by requiring one measure to be
consistent across all scales. They occur infinitely often for a single
compact planar set with the prescribed exact dimensions. Conversely, the
single-collapse profile cannot occur at every large scale. Neither statement
is a counterexample to a distance-set theorem.

\begin{theorem}[Subsequential realization]\label{realization:theorem}
Let \(0<a<b\le1\), \(d=1+a\), \(D=1+b\). Suppose that
\(g:[0,1]\to\R\) is 1-Lipschitz, satisfies
\(ax\le g(x)\le bx\), has \(g(1)=a\), and equals \(bx\) on
some interval \([0,c]\), \(c>0\). There are a compact set
\(E\subset[0,1]^2\), a \(d\)-Frostman probability \(\mu\) on \(E\),
and integers \(N_j\to\infty\), with
\[
 \HH E=d,\qquad \PP E=\overline{\dim}_{\rm B}E=D,
\]
such that every occupied half-open dyadic square of depth \(n\) has mass
\(2^{-F_{\rm dyad}(n)}\). With linear interpolation of \(F_{\rm dyad}\),
\begin{equation}\label{realization:convergence}
 \sup_{x\in[0,1]}
 \left|\frac{F_{\rm dyad}(N_jx)}{N_j}-x-g(x)\right|\longrightarrow0.
\end{equation}
The associated limiting optimal chain costs converge to \(\Phi_0(g)\).
\end{theorem}

\begin{proof}
Put \(h(x)=x+g(x)\). It is nondecreasing and 2-Lipschitz, with
\(dx\le h(x)\le Dx\), \(h(1)=d\), and \(h(x)=Dx\) for
\(x\le c\). Choose integers \(N_0=0<N_1<\cdots\) with
\(N_{j-1}/N_j<c\) and \(N_{j-1}/N_j\to0\). On
\([N_{j-1},N_j]\), define
\begin{equation}\label{realization:construction}
 F(t)=\max\{dt,\ N_jh(t/N_j)-(D-d)N_{j-1}\}.
\end{equation}
At the left endpoint both entries equal \(dN_{j-1}\), and at the
right endpoint the maximum equals \(dN_j\). Thus \(F\) is continuous,
nondecreasing, and 2-Lipschitz, and
\[
 dt\le F(t)\le Dt,\qquad F(N_j)=dN_j.
\]
For every \(0\le t\le N_j\),
\begin{equation}\label{realization:approximation}
 0\le N_jh(t/N_j)-F(t)\le(D-d)N_{j-1}.
\end{equation}
For \(t\ge N_{j-1}\), this follows from the defining maximum; for
\(t\le N_{j-1}\), use \(h(t/N_j)=Dt/N_j\) and \(dt\le F(t)\le Dt\).
At a fixed \(c_0\in(0,c]\), the ratio \(F(c_0N_j)/(c_0N_j)\)
therefore tends to \(D\). It follows that
\[
 \liminf_{t\to\infty}F(t)/t=d,\qquad
 \limsup_{t\to\infty}F(t)/t=D.
\]

Set \(S(n)=\lfloor F(n)/2\rfloor\). The increments of \(S\) are
zero or one. Let \(A\) be the set of indices with increment one, and put
\[
 C_A=\left\{\sum_{n\ge1}\varepsilon_n2^{-n}:
 \varepsilon_n\in\{0,1\}\ (n\in A),\quad
 \varepsilon_n=0\ (n\notin A)\right\},\qquad E=C_A\times C_A.
\]
Give every allowed binary digit in each coordinate an independent fair
choice. The resulting probability \(\mu\) gives each occupied depth-\(n\)
square mass \(2^{-2S(n)}\), and there are \(2^{2S(n)}\) such squares.
Ambiguous binary expansions have zero measure, because there are infinitely
many free digits and a prescribed infinite tail has probability zero.
The support is compact. With \(F_{\rm dyad}(n)=2S(n)\), rounding and
interpolation give
\[
 \|F_{\rm dyad}-F\|_{L^\infty([0,\infty))}\le4.
\]
Equation~\eqref{realization:approximation} now proves
\eqref{realization:convergence} with error at most
\(((D-d)N_{j-1}+4)/N_j\).

A ball of radius \(2^{-n}\) meets a bounded number of depth-\(n\)
dyadic squares. Consequently
\begin{equation}\label{realization:ball}
 \mu(B(z,2^{-n}))\le C2^{-F_{\rm dyad}(n)}\le4C2^{-dn}.
\end{equation}
Neighboring dyadic radii give the \(d\)-Frostman bound at all radii,
so \(\HH E\ge d\). At \(n=N_j\), the support is covered by at most
\(2^{dN_j}\) squares of diameter \(\sqrt2\,2^{-N_j}\). Their
\(s\)-dimensional covering sum tends to zero for every \(s>d\), proving
\(\HH E=d\).

The occupied-square count and the first inequality in
\eqref{realization:ball} give matching upper and lower covering-number
bounds, within absolute constants. Hence
\(\overline{\dim}_{\rm B}E=\limsup_nF_{\rm dyad}(n)/n=D\).
For the packing dimension, use the standard characterization
\begin{equation}\label{realization:packing}
 \PP E=\inf_{E\subset\bigcup_iE_i}
       \sup_i\overline{\dim}_{\rm B}E_i,
\end{equation}
where the covers are countable and bounded; see, for example,
\href{https://research-repository.st-andrews.ac.uk/bitstream/handle/10023/17263/CapacityProjections_Falconer.pdf?sequence=1}
{Falconer, \emph{A capacity approach to box and packing dimensions},
equation (4.1)}.
Some member of any such cover has \(\mu^*(E_i\cap E)>0\).
If \(M\) balls of radius \(2^{-n}\) cover that member, then
\[
 \mu^*(E_i\cap E)\le MC2^{-F_{\rm dyad}(n)}.
\]
Thus \(\overline{\dim}_{\rm B}E_i\ge D\). Equation~\eqref{realization:packing}
proves \(\PP E\ge D\), and the upper box bound gives equality.
The same argument shows that every subset of positive \(\mu\) outer measure
has packing dimension \(D\).

It remains to justify convergence of the costs; this follows from the
stability lemma below. The interpolated functions
\(F_{\rm dyad}(Nx)/N-x\) are 1-Lipschitz, since their slopes are
\(-1\) or \(1\).
\end{proof}

\begin{lemma}[Continuity of limiting cost]\label{realization:stability}
For 1-Lipschitz functions \(g,h\) on \([0,1]\), let \(\Phi_0\) denote
the limiting optimal cost of finite admissible chains as the starting point
tends to one. If \(\|g-h\|_\infty\le\epsilon\), then, for
\(0<\zeta<1\),
\[
 |\Phi_0(g)-\Phi_0(h)|\le2\zeta+2K(\zeta)\epsilon,
 \qquad K(\zeta)=2\lceil\log_2(1/\zeta)\rceil+4.
\]
In particular \(\Phi_0\) is continuous under uniform convergence within
this Lipschitz class.
\end{lemma}
\begin{proof}
Write \(\Psi_g(n)\) for the infimum of costs from \(n<1\) to zero.
If \(n<n'\), truncate a chain from \(n'\) at its first crossing of \(n\).
The crossing edge remains admissible and its cost can only decrease. Conversely,
extend a chain from \(n\) up to \(n'\) by repeatedly doubling the
complementary distance \(1-x\), stopping at \(n\). The added cost is
at most the negative variation on \([n,n']\), hence at most \(n'-n\).
Thus
\[
 0\le\Psi_g(n')-\Psi_g(n)\le n'-n,
 \qquad 0\le\Phi_0(g)-\Psi_g(1-\zeta)\le\zeta.
\]
Merging adjacent admissible edges never increases cost, since
\(c_g(l,n)\le c_g(l,m)+c_g(m,n)\) for \(l<m<n\).
After all possible mergers, every two edges more than double \(1-x\),
giving the bound \(K(\zeta)\). Each edge cost changes by at most
\(2\epsilon\) under the uniform perturbation. Infimizing over the
bounded-length chains and adding the two truncation errors proves the claim.
\end{proof}

\begin{remark}[Applicable templates and separated restrictions]
\label{realization:templates}
The single-collapse profile \(\min\{bx,1+a-x\}\) satisfies the realization
hypotheses for every \(0<a<b\le1\). The matching two-collapse profile
\(\min\{bx,at+|x-t|,1+a-x\}\), in its certified parameter region,
also does: its two other entries are strictly positive at zero, so it equals
\(bx\) on an initial interval. Therefore both sharp costs occur along
subsequences for actual Frostman measures.

A fixed occupied depth-\(k\) cylinder, with the measure normalized on it,
has descendant cube exponent \(F_{\rm dyad}(n)-F_{\rm dyad}(k)\).
The bounded initial change disappears after normalization by \(N_j\), so
it has the same limiting profiles and dimensions. Choosing two sufficiently
small cylinders around distinct support points gives compact source and pin
supports a positive distance apart, each with these properties.
\end{remark}

\begin{lemma}[Persistence under large regular restrictions]
\label{realization:regularization}
For the measure in Theorem~\ref{realization:theorem}, fix a union \(X\)
of depth-\(N\) cubes, with \(w=\mu(X)>0\), and let
\(\sigma=\mu|_X/w\). Suppose its occupied depth-\(k\) cubes obey
\[
 2^{-f(k)-eN}\le\sigma(Q)\le2^{-f(k)}.
\]
Then
\begin{equation}\label{realization:sandwich}
 F_{\rm dyad}(k)+\log_2w-eN\le f(k)\le F_{\rm dyad}(k).
\end{equation}
In particular, a component with \(w\ge2^{-\gamma N}\) changes the
profile by at most \((\gamma+e)N\).
\end{lemma}
\begin{proof}
For the lower bound, combine the lower regularity estimate with
\(\sigma(Q)\le\mu(Q)/w=2^{-F_{\rm dyad}(k)}/w\).
For the upper bound, at most \(2^{F_{\rm dyad}(k)}\) occupied cubes
have masses at most \(2^{-f(k)}\), and their masses sum to one.
\end{proof}

Thus selecting large regular components cannot remove the matching profiles
at the constructed scales when the regularization losses tend to zero.
This does not preclude using geometric information absent from the mass profile.

\subsection{Single collapse cannot persist at all scales}
\label{realization:sparsity-section}
Let \(G:[0,\infty)\to\R\) be continuous and put
\(g_N(x)=G(Nx)/N\). Take
\(g(x)=\min\{bx,1+a-x\}\), \(0<a<b\le1\).

\begin{proposition}[Separated clusters of near occurrences]
\label{realization:sparsity}
Suppose \(0<M<N\) and both \(g_M,g_N\) are within uniform distance
\(\epsilon\) of \(g\), where \(0<\epsilon<b-a\). Then
\[
 \frac MN\le L_\epsilon:=\frac{\epsilon}{b-a-\epsilon}
 \quad\text{or}\quad
 \frac MN\ge U_\epsilon:=\frac{1+a-\epsilon}{1+a+\epsilon}.
\]
For sufficiently small \(\epsilon>0\), write
\(\ell_\epsilon=\log(1/L_\epsilon)\) and
\(w_\epsilon=\log(1/U_\epsilon)\). Then
\begin{equation}\label{realization:log-density}
 \limsup_{L\to\infty}\frac1L
 \left|\left\{t\in[0,L]:\|g_{e^t}-g\|_\infty\le\epsilon\right\}\right|
 \le\frac{w_\epsilon}{\ell_\epsilon}
 =O_{a,b}\!\left(\frac{\epsilon}{\log(1/\epsilon)}\right).
\end{equation}
\end{proposition}
\begin{proof}
Put \(r=M/N\) and \(p=(1+a)/(1+b)\). The endpoint of \(g_M\)
and the value of \(g_N\) at \(r\) give
\[
 |g(r)-ra|\le\epsilon(1+r).
\]
For \(r\le p\), the left side equals \(r(b-a)\), yielding
\(r\le L_\epsilon\). For \(r\ge p\), it equals
\((1+a)(1-r)\), yielding \(r\ge U_\epsilon\).
For small \(\epsilon\), one has \(0<L_\epsilon<U_\epsilon<1\).
Thus in logarithmic scale no two near occurrences have separation strictly
between \(w_\epsilon\) and \(\ell_\epsilon\). In an interval of
length less than \(\ell_\epsilon\), their set has diameter and measure
at most \(w_\epsilon\). It is measurable by continuity of \(G\).
Partition long intervals into pieces of length
\(\ell_\epsilon-\delta\), and then let \(\delta\downarrow0\), to get
\eqref{realization:log-density}. Finally,
\(w_\epsilon=O_a(\epsilon)\) and
\(\ell_\epsilon\sim\log(1/\epsilon)\).
\end{proof}

For either matching template, a simpler obstruction follows from its initial
linear interval: if \(\lambda\le c\) and both scales \(N,\lambda N\)
have error at most \(\epsilon\), then
\(\lambda(b-a)\le(1+\lambda)\epsilon\).

These results distinguish infinite recurrence from recurrence at every scale.
The sparsity bound does not make shell estimates of size \(2^{cN_j}\)
summable when \(c>0\). Moreover, profiles obtained by separately regularizing
an arbitrary measure at each scale need not be prefixes of a single function
\(G\). A new argument would be needed to exploit this obstruction in the
analytic transfer. No improvement to the distance criterion is claimed here.

\subsection{Simultaneous raw Fourier-energy saturation}
\label{realization:fourier-section}

The profiles can also occur where the source's raw Fourier energy nearly
attains the Frostman exponent. The assertion concerns fixed separated
restrictions of the same measure, so it is compatible with the source--pin
separation required in the analytic proof.

\begin{proposition}[Near-extremal profiles and energy on one annulus]
\label{realization:fourier}
Under Theorem~\ref{realization:theorem}, choose two separated occupied
cylinders at a fixed common depth \(k\), and normalize the restrictions
of \(\mu\) to obtain source and pin probabilities \(\sigma,\nu\).
Their compact supports both have Hausdorff dimension \(d\) and packing
dimension \(D\), and both probabilities are \(d\)-Frostman.
Writing \(F=F_{\rm dyad}\), they have the common conditional exponent
\[
 F_k(t)=\begin{cases}0,&0\le t\le k,\\
 F(t)-F(k),&t\ge k.
 \end{cases}
\]
For every \(\tau>0\) and \(\epsilon>0\), there are infinitely many
integers \(n\) such that
\begin{equation}\label{realization:fourier-conclusion}
 \left\|F_k(nx)/n-x-g(x)\right\|_{L^\infty_x([0,1])}\le\tau,
 \qquad
 \int_{2^n\le|\xi|<2^{n+1}}|\widehat\sigma(\xi)|^2\dd\xi
 \ge2^{(2-d-\epsilon)n}.
\end{equation}
The source and pin supports are fixed independently of both tolerances.
\end{proposition}

\begin{proof}
The restrictions and dimensions were justified in
Remark~\ref{realization:templates}. They can be chosen at the same depth
by taking two sufficiently small dyadic squares with separated closures.
Both original cube masses are \(2^{-F(k)}\). Every occupied descendant
at depth \(n\ge k\) has conditional mass \(2^{F(k)-F(n)}\),
proving the formula for \(F_k\). In particular
\begin{equation}\label{realization:conditional-error}
 \|F_k-F\|_\infty\le F(k)\le2k.
\end{equation}

Use the convention \(\widehat\sigma(\xi)=\int e^{-2\pi i x\cdot\xi}
\dd\sigma(x)\). The Gaussian Fourier transform and Fubini give
\begin{align}
 \mathcal E_\sigma(R)
 &:=\int e^{-\pi|\xi|^2/R^2}|\widehat\sigma(\xi)|^2\dd\xi
 \nonumber\\
 &=R^2\iint e^{-\pi R^2|x-y|^2}\dd\sigma(x)\dd\sigma(y).
 \label{realization:gaussian}
\end{align}
Fubini is justified by Gaussian integrability and finiteness of the measure.
For \(R_j=2^{N_j}\), two points in the same depth-\(N_j\) square
have distance at most \(\sqrt2/R_j\). The sum of the squares of the
conditional masses of these squares is \(2^{F(k)-F(N_j)}\).
Since \(F(N_j)\le dN_j\),
\begin{equation}\label{realization:energy-lower}
 \mathcal E_\sigma(R_j)\ge e^{-2\pi}R_j^2 2^{F(k)-F(N_j)}
 \ge c_0R_j^{2-d},\qquad c_0=e^{-2\pi}2^{F(k)}>0.
\end{equation}

Conversely, the Frostman bound and spatial annuli of radii \(2^m/R\)
give
\[
 \int e^{-\pi R^2|x-y|^2}\dd\sigma(y)
 \le CR^{-d}\left(1+\sum_{m\ge1}2^{md}e^{-\pi2^{2m-2}}\right).
\]
The sum is finite. Equation~\eqref{realization:gaussian} therefore yields,
for \(R\ge1\),
\begin{equation}\label{realization:energy-upper}
 \mathcal E_\sigma(R)\le C_1R^{2-d},\qquad
 \int_{|\xi|\le R}|\widehat\sigma(\xi)|^2\dd\xi
 \le e^\pi\mathcal E_\sigma(R)\le C_2R^{2-d}.
\end{equation}

Put \(q=2-d>0\), and temporarily fix \(0<\eta<1/4\).
At \(R=R_j\), the part of the Gaussian integral below
\(R_j^{1-\eta}\) is at most \(C_2R_j^{(1-\eta)q}=o(R_j^q)\).
The part above \(R_j^{1+\eta}\), using
\(|\widehat\sigma|\le1\), is at most the exact planar Gaussian tail
\[
 \int_{|\xi|\ge R_j^{1+\eta}}e^{-\pi|\xi|^2/R_j^2}\dd\xi
 =R_j^2e^{-\pi R_j^{2\eta}}=o(R_j^q).
\]
Hence the intermediate region contributes at least
\((c_0/2)R_j^q\) for large \(j\). Its Gaussian weight is at most one,
and it is covered by at most \(2\eta N_j+3\) dyadic annuli.
Some index \(n_j\) consequently satisfies
\begin{gather}
 (1-\eta)N_j-1\le n_j\le(1+\eta)N_j+1,
 \label{realization:nearby-annulus}\\
 \int_{2^{n_j}\le|\xi|<2^{n_j+1}}|\widehat\sigma(\xi)|^2\dd\xi
 \ge\frac{c_\eta}{N_j}2^{qN_j},\qquad c_\eta>0.
 \label{realization:annular-lower}
\end{gather}
Only existence of one annulus in this window is asserted.

For the given tolerances, take
\[
 \epsilon_0=\min\{\epsilon,q/2\},\qquad
 0<\eta<\min\{1/4,\tau/16,\epsilon_0/(8q)\}.
\]
For large \(j\),
\[
 (q-\epsilon_0)n_j
 \le(q-\epsilon_0)(1+2\eta)N_j
 \le(q-3\epsilon_0/4)N_j.
\]
The polynomial loss in \eqref{realization:annular-lower} is absorbed by
this strict exponential margin, proving the energy inequality in
\eqref{realization:fourier-conclusion}.

Finally \(G(t)=F(t)-t\) is 1-Lipschitz and \(G(0)=0\), so
\[
 \left\|G(nx)/n-G(Nx)/N\right\|_{L^\infty_x([0,1])}
 \le2|n-N|/n.
\]
This follows by separating the changes of argument and denominator, using
\(|G(Nx)|\le N\). Equations~\eqref{realization:conditional-error}
and \eqref{realization:nearby-annulus} give, for large \(j\),
\[
 \|F_k(n_jx)/n_j-x-g(x)\|_\infty
 \le F(k)/n_j+4\eta+4/N_j+
       \|G(N_jx)/N_j-g(x)\|_\infty.
\]
The second term is less than \(\tau/4\), and each of the other terms
is at most \(\tau/4\) for sufficiently large \(j\), by
\eqref{realization:convergence}. This proves the profile inequality.
The windows in \eqref{realization:nearby-annulus} are eventually disjoint,
since \(N_{j-1}/N_j\to0\), so infinitely many distinct indices occur.
\end{proof}

This proposition rules out a universal saving in the \emph{raw} source
Fourier-energy exponent merely from the pin profile being near one of the
matching templates. It proves no sharpness for the modified good-packet
estimate and says nothing adverse about the actual pinned distance sets.
Angular geometry, cancellation, or a different summation argument could
still yield a stronger distance theorem.
